\documentclass[12pt]{article}
\usepackage{amsmath, amssymb}
\usepackage[margin=1in]{geometry}
\setlength{\parskip}{6pt}

\title{QUBO Formulation: N-Queens Exact-Cover Problem}
\date{}

\begin{document}
\maketitle

\subsection*{N-Queens Exact-Cover Problem}

\paragraph{Problem Statement.}
Given an $N \times N$ chessboard, place $N$ queens such that no two queens attack each other. A queen at position $(i,j)$ attacks all squares in row $i$, column $j$, and both diagonals passing through $(i,j)$. The task is to select exactly $N$ squares satisfying these non-attacking conditions.

\paragraph{Decision Variables.}
We introduce a binary variable $x_{i,j} \in \{0,1\}$ for each square $(i,j)$ on the board, where $i,j \in \{0,1,\dots,N-1\}$. The variable $x_{i,j} = 1$ if and only if a queen is placed at row $i$ and column $j$; otherwise $x_{i,j} = 0$.

\paragraph{QUBO Derivation.}
The problem is a feasibility (exact-cover) problem with no intrinsic cost function. We formulate it as minimizing the total number of violated constraints by converting each constraint into a quadratic penalty term. Let $\lambda > 0$ be the penalty weight.

\textbf{Step 1 -- Objective.}
Since there is no original objective function, we minimize the sum of squared violations of all constraints:
\begin{align}
\min_{x \in \{0,1\}^{N^2}} \left[ \text{Row Violations} + \text{Column Violations} + \text{Diagonal Violations} + \text{Anti-Diagonal Violations} \right].
\end{align}

\textbf{Step 2 -- Constraint Algebra.}
The constraints are expressed as follows:
\begin{align}
\text{Row } i: & \quad \sum_{j=0}^{N-1} x_{i,j} = 1, \quad \forall i \in \{0,\dots,N-1\}, \\
\text{Column } j: & \quad \sum_{i=0}^{N-1} x_{i,j} = 1, \quad \forall j \in \{0,\dots,N-1\}, \\
\text{Main Diagonal } k: & \quad \sum_{(i,j) \in D_k} x_{i,j} \le 1, \quad \forall k \in \{-(N-1),\dots,N-1\}, \\
\text{Anti-Diagonal } k: & \quad \sum_{(i,j) \in AD_k} x_{i,j} \le 1, \quad \forall k \in \{0,\dots,2N-2\},
\end{align}
where $D_k = \{(i,j) : i-j=k\}$ and $AD_k = \{(i,j) : i+j=k\}$.

\textbf{Step 3 -- Penalty Expansion.}
We convert each constraint to a quadratic penalty. For equality constraints (rows and columns), we use the squared form $\lambda(\sum x - 1)^2$. For inequality constraints (diagonals), we penalize pairwise conflicts using $\lambda \sum_{u<v} x_u x_v$, which is zero when at most one variable is active and positive otherwise.

\textit{Row constraints:}
\begin{align}
P_{\text{row}}(i) &= \lambda \left( \sum_{j=0}^{N-1} x_{i,j} - 1 \right)^2 \\
&= \lambda \left( \sum_{j=0}^{N-1} x_{i,j}^2 - 2 \sum_{0 \le j < l \le N-1} x_{i,j}x_{i,l} + 1 \right) \\
&= \lambda \sum_{j=0}^{N-1} x_{i,j} - 2\lambda \sum_{0 \le j < l \le N-1} x_{i,j}x_{i,l} + \lambda,
\end{align}
where we used the idempotent property $x_{i,j}^2 = x_{i,j}$.

\textit{Column constraints:}
\begin{align}
P_{\text{col}}(j) &= \lambda \left( \sum_{i=0}^{N-1} x_{i,j} - 1 \right)^2 \\
&= \lambda \sum_{i=0}^{N-1} x_{i,j} - 2\lambda \sum_{0 \le i < i' \le N-1} x_{i,j}x_{i',j} + \lambda.
\end{align}

\textit{Diagonal constraints:}
\begin{align}
P_{\text{diag}}(k) &= \lambda \sum_{(i,j) < (i',j') \in D_k} x_{i,j}x_{i',j'}, \\
P_{\text{adiag}}(k) &= \lambda \sum_{(i,j) < (i',j') \in AD_k} x_{i,j}x_{i',j'}.
\end{align}

\textbf{Step 4 -- Combined Hamiltonian.}
Summing all penalty contributions yields the QUBO Hamiltonian $H(x)$:
\begin{align}
H(x) &= \sum_{i=0}^{N-1} P_{\text{row}}(i) + \sum_{j=0}^{N-1} P_{\text{col}}(j) + \sum_{k} P_{\text{diag}}(k) + \sum_{k} P_{\text{adiag}}(k).
\end{align}

\textbf{Step 5 -- Q Matrix and Offset.}
Collecting linear and quadratic terms from the expanded penalties, we read off the QUBO coefficients $Q_{u,v}$ and constant offset $c$ such that $H(x) = \sum_{u \le v} Q_{u,v} x_u x_v + c$. Identifying each board square $(i,j)$ with a unique index $u$, we obtain:
\begin{align}
Q_{(i,j),(i,j)} &= 2\lambda, \\
Q_{(i,j),(i,l)} &= -2\lambda \quad \text{for } j \neq l, \\
Q_{(i,j),(i',j)} &= -2\lambda \quad \text{for } i \neq i', \\
Q_{(i,j),(i',j')} &= \lambda \quad \text{if } (i,j) \text{ and } (i',j') \text{ share a main or anti-diagonal}, \\
Q_{(i,j),(i',j')} &= 0 \quad \text{otherwise}, \\
c &= 2N\lambda.
\end{align}

\paragraph{Penalty Weight Justification.}
The penalty weight $\lambda$ must be large enough to ensure that any energy reduction from satisfying a constraint outweighs the penalty incurred by violating another. Since the problem contains no intrinsic objective terms (it is purely a feasibility problem), there is no trade-off between objective improvement and constraint violation. Therefore, any $\lambda > 0$ strictly penalizes infeasible configurations relative to feasible ones. We set $\lambda = 1$ for simplicity, which matches the standard scaling for exact-cover QUBOs.

\paragraph{Correctness Argument.}
Minimizing $H(x)$ over $\{0,1\}^{N^2}$ recovers the optimal solution because: (i) any feasible placement satisfies all constraints exactly, yielding zero penalty terms and thus $H(x) = 0$; (ii) any infeasible placement violates at least one constraint, resulting in strictly positive penalty terms and thus $H(x) > 0$. Consequently, the global minima of $H(x)$ are exactly the feasible queen placements, and the minimum energy is $0$.

\paragraph{Illustrative Example.}
Consider the concrete instance $N=4$. The board has $4 \times 4 = 16$ squares, indexed $0$ to $15$ via $u = 4i + j$. The decision variables are $x_{0,0}, \dots, x_{3,3} \in \{0,1\}$.

Using $\lambda = 1$, the general formulas instantiate as follows:
\begin{itemize}
    \item \textbf{Row penalties:} For each $i \in \{0,1,2,3\}$, $P_{\text{row}}(i) = \sum_{j=0}^{3} x_{i,j} - 2\sum_{0 \le j < l \le 3} x_{i,j}x_{i,l} + 1$.
    \item \textbf{Column penalties:} For each $j \in \{0,1,2,3\}$, $P_{\text{col}}(j) = \sum_{i=0}^{3} x_{i,j} - 2\sum_{0 \le i < i' \le 3} x_{i,j}x_{i',j} + 1$.
    \item \textbf{Diagonal penalties:} There are $7$ main diagonals and $7$ anti-diagonals. Each contributes pairwise terms with coefficient $1$. For example, the main diagonal $k=0$ contains squares $(0,0), (1,1), (2,2), (3,3)$, yielding $x_{0,0}x_{1,1} + x_{0,0}x_{2,2} + x_{0,0}x_{3,3} + x_{1,1}x_{2,2} + x_{1,1}x_{3,3} + x_{2,2}x_{3,3}$.
\end{itemize}

The resulting Hamiltonian is $H(x) = \sum_{u \le v} Q_{u,v} x_u x_v + 8$. The constant offset is $c = 2(4)(1) = 8$.

A feasible solution is given by queens at coordinates $(0,1), (1,3), (2,0), (3,2)$, corresponding to $x_{0,1}=x_{1,3}=x_{2,0}=x_{3,2}=1$ and all other variables $0$. Substituting this assignment into $H(x)$ yields $H(x) = 0$, confirming feasibility and optimality. Any assignment with two queens sharing a row, column, or diagonal will produce $H(x) \ge 1$.

\end{document}